% !TeX spellcheck = pl_PL-Polish
\documentclass[12pt,a4paper]{article}

\usepackage[T1]{fontenc}
\usepackage[polish]{babel}
\usepackage[utf8]{inputenc}
\usepackage{lmodern}
\selectlanguage{polish}
\usepackage{indentfirst}
\usepackage{graphicx}
\usepackage{xcolor}
\usepackage{hyperref}
\usepackage{float}
\usepackage{booktabs}
\usepackage{array}
\usepackage{longtable}
\usepackage{enumitem}
\usepackage{listings}
\usepackage{microtype}

\def\projectName{SCHIZA -- System Ciągłej Hybrydowej Identyfikacji Zagrożeń i Anomalii}
\def\authorA{Dawid Pągowski}
\def\authorB{Bartłomiej Łokaj}
\def\authorC{Paweł Świderski}

\newcommand\todo[1]{\textcolor{red}{#1}}
\newcommand{\code}[1]{\texttt{#1}}
\newcommand{\project}{Schiza}

\definecolor{codebackground}{RGB}{245,245,245}
\definecolor{codecomment}{RGB}{80,100,80}
\definecolor{linkcolor}{RGB}{0,70,130}

\hypersetup{
    colorlinks=true,
    linkcolor=linkcolor,
    urlcolor=linkcolor,
    pdftitle={SCHIZA -- System Ciągłej Hybrydowej Identyfikacji Zagrożeń i Anomalii},
    pdfauthor={Zespół projektowy}
}

\lstset{
    basicstyle=\ttfamily\small,
    backgroundcolor=\color{codebackground},
    frame=single,
    breaklines=true,
    columns=fullflexible,
    keepspaces=true,
    showstringspaces=false,
    commentstyle=\color{codecomment}
}

\begin{document}
\pagenumbering{gobble}
\clearpage
\begin{figure}[h]
    \centering
    {\includegraphics[width=0.5\linewidth]{../prezentacje/media/logo.jpg}}{\vspace{2cm}}
\end{figure}

\hspace{3cm}
\begin{center}Dokumentacja projektowa\end{center}
\hspace{3cm}
\begin{center}\large\textbf{Inżynieria Oprogramowania}\end{center}
\begin{center}\large\textit{\projectName}\end{center}

\hspace{7cm}
\begin{flushright}
    Kierunek: Informatyka \par
    \textit{\authorA} \par
    \textit{\authorB} \par
    \textit{\authorC} \par
\end{flushright}
\vfill
\begin{center}Gliwice, 2026\end{center}

\newpage
\pagenumbering{arabic}
\tableofcontents
\newpage

\section{Cel i zakres projektu}

\project{} jest skrótem od nazwy \textbf{System Ciągłej Hybrydowej Identyfikacji
Zagrożeń i Anomalii}. Nazwa odzwierciedla podstawowe założenia i rzeczywisty
zakres implementacji projektu: monitorowanie kolejnych wpisów logów serwera
Apache, łączenie kilku metod analizy oraz identyfikowanie żądań HTTP
wskazujących na potencjalne zagrożenia lub nietypowe zachowania.

Poszczególne człony nazwy mają następujące znaczenie:
\begin{itemize}
    \item \textbf{System} -- rozwiązanie obejmujące parsery logów, backend
    w C\#, moduł detekcji w Pythonie oraz frontend webowy;
    \item \textbf{Ciągłej} -- backend cyklicznie odczytuje nowe wpisy z pliku
    logu Apache i przekazuje poprawnie sparsowane żądania do dalszej analizy;
    \item \textbf{Hybrydowej} -- moduł detekcji łączy analizę opartą na
    deterministycznych regułach bezpieczeństwa z analizą wykorzystującą model
    uczenia maszynowego typu autoencoder;
    \item \textbf{Identyfikacji} -- system analizuje ustandaryzowane dane
    żądań HTTP i klasyfikuje je jako podejrzane lub nieoznaczone jako anomalne;
    \item \textbf{Zagrożeń} -- reguły detekcji obejmują charakterystyczne
    wzorce związane między innymi z SQL injection, XSS, path traversal,
    ujawnieniem plików i command injection;
    \item \textbf{Anomalii} -- model analizuje wektor cech żądania, a moduł
    detekcji zwraca wynik punktowy, decyzję oraz informacje o poziomie
    istotności i powodach klasyfikacji.
\end{itemize}

Celem projektu jest zbudowanie systemu, który obserwuje logi serwera Apache,
parsuje żądania HTTP, przekazuje ich ustandaryzowaną reprezentację do modułu
wykrywania anomalii, a następnie zapisuje i prezentuje zdarzenia uznane za
podejrzane.

System został podzielony na trzy główne obszary:
\begin{enumerate}
    \item parsery logów oraz współdzielone modele żądań,
    \item backend odpowiedzialny za przetwarzanie, komunikację, zapis i API,
    \item moduł detekcji Python oraz aplikacja webowa prezentująca wyniki.
\end{enumerate}

Głównym scenariuszem jest ciągłe odczytywanie nowych linii z pliku
\code{access.log}. Dla każdej poprawnie sparsowanej linii backend buduje obiekt
\code{ApacheRequest}, wysyła wybrane dane do detektora Python i zapisuje pełne
zdarzenie, jeżeli detektor zwróci \code{is\_anomaly = true}.

\section{Organizacja i podział odpowiedzialności}

Podział odpowiedzialności przyjęty w projekcie przedstawiono w tabeli
\ref{tab:roles}. Nazwy osób pozostawiono w formie ról zgodnie z dostarczonym
podziałem zespołu.

\begin{table}[H]
    \centering
    \caption{Podział odpowiedzialności w projekcie}
    \label{tab:roles}
    \begin{tabular}{p{0.18\textwidth}p{0.72\textwidth}}
        \toprule
        \textbf{Rola} & \textbf{Zakres odpowiedzialności} \\
        \midrule
        \authorA & Lead projektu oraz właściciel repozytorium i prototypu projektu
        na GitHubie. Autor parsera logów w katalogu \code{/logEngine},
        komponentów współdzielonych w \code{/shared} oraz mocków
        wykorzystywanych do testowania frontendu i parsera. \\
        \authorB & Autor modułu detekcji opartego na Pythonie
        (\code{/Deep Learning Detection Module}) oraz backendu
        zaimplementowanego w C\#. \\
        \authorC & Autor frontendowej części aplikacji. \\
        \bottomrule
    \end{tabular}
\end{table}


\section{Architektura systemu}

\subsection{Podział na projekty i warstwy}

Rozwiązanie .NET jest zapisane w pliku \code{Schiza.slnx} i składa się z
następujących projektów:

\begin{itemize}
    \item \code{backend/Schiza.Domain} -- encje domenowe, w szczególności
    \code{FlaggedEvent},
    \item \code{backend/Schiza.Application} -- porty, przypadki użycia
    i logika aplikacyjna niezależna od infrastruktury,
    \item \code{backend/Schiza.Infrastructure} -- SQLite, worker logów,
    klient HTTP detektora oraz implementacje infrastrukturalne,
    \item \code{backend/Schiza.Api} -- aplikacja ASP.NET Core, kontrolery,
    uwierzytelnianie i rejestracja zależności,
    \item \code{backend/Schiza.Shared} -- wspólne modele żądań i konwertery JSON,
    \item \code{backend/Schiza.LogEngine} -- niezależny parser logów Apache
    i parser audit logów ModSecurity.
\end{itemize}

Zależności warstw można przedstawić następująco:

\begin{verbatim}
API -> Application -> Domain
 |         |
 +-> Infrastructure -> LogEngine, Shared, SQLite, Python HTTP API

Shared jest wykorzystywany przez Application, Infrastructure i LogEngine.
\end{verbatim}

Kontrolery nie komunikują się bezpośrednio z SQLite ani z modułem Python.
Wykorzystują interfejsy warstwy Application, a ich implementacje są dostarczane
przez Infrastructure.

\subsection{Główny przepływ danych}

\begin{verbatim}
Apache access.log
       |
       v
ApacheLogWorker -- odczyt nowych linii co około 2 sekundy
       |
       v
ApacheLogParser
       |
       +-- Unparseable -> ostrzeżenie w logu backendu
       |
       v
ApacheRequest
       |
       v
LogEventProcessor
       |
       v
PythonAnomalyDetector -- POST /predict
       |
       +-- is_anomaly = false -> brak zapisu
       |
       +-- is_anomaly = true -> SQLite FlaggedEvents
                                      |
                                      v
                              API + frontend React
\end{verbatim}

\section{Przetwarzanie logów Apache}

\subsection{Format wejściowy}

Parser \code{ApacheLogParser} obsługuje rozszerzony format Apache odpowiadający
następującej dyrektywie \code{LogFormat}:

\begin{lstlisting}
LogFormat "%h %l %u %t \"%r\" %>s %O \"%{Referer}i\" \"%{User-Agent}i\" %D \"%{Host}i\" %v:%p %{SSL_PROTOCOL}x %B %{UNIQUE_ID}e" combined_ext
\end{lstlisting}

Parser rozpoznaje między innymi:
\begin{itemize}
    \item adres IP klienta i timestamp,
    \item request line, w tym metodę, ścieżkę, query string i protokół,
    \item status HTTP, liczbę wysłanych bajtów i czas trwania żądania,
    \item referer, User-Agent, host i protokół TLS,
    \item opcjonalny rozmiar body oraz wartość \code{UNIQUE\_ID}.
\end{itemize}

Wartość \code{UNIQUE\_ID} może zostać wykorzystana do powiązania wpisu
\code{access.log} z odpowiednim wpisem audit logu ModSecurity.

\subsection{Odporność parsera}

Parser nie rzuca wyjątku dla typowych błędów formatu. Zwraca wynik typu
\code{LogParseResult}, który może być wariantem \code{Parsed} albo
\code{Unparseable} z kodem błędu, na przykład \code{Empty},
\code{UnexpectedStructure}, \code{InvalidClientIp}, \code{InvalidTimestamp}
lub \code{InvalidNumber}.

Niekompletna linia logu nie jest przekazywana do detektora. Jest rejestrowana
jako ostrzeżenie przez \code{ApacheLogWorker}. W środowisku produkcyjnym należy
uwzględnić, że odczyt pliku w trakcie zapisu może wymagać dodatkowego mechanizmu
buforowania niepełnych linii.

\section{Modele współdzielone i ModSecurity}

\subsection{ApacheRequest}

\code{ApacheRequest} jest głównym modelem żądania używanym w backendzie. Zawiera:
\begin{itemize}
    \item \code{ClientIp}, \code{Timestamp} i \code{RequestLine},
    \item referer, User-Agent, host i protokół TLS,
    \item \code{UniqueId} z modułu \code{mod\_unique\_id},
    \item opcjonalne dane odpowiedzi \code{ResponseInfo},
    \item opcjonalne \code{RequestBody} zawierające typ treści i body żądania.
\end{itemize}

Wartości adresów IP są serializowane jako tekst. Nazwy JSON są zgodne z opcjami
\code{JsonSerializerOptions.Web}, co pozwala bezpośrednio mapować dane na kontrakty
innych komponentów.

\subsection{Parser audit logu ModSecurity}

Moduł \code{Schiza.LogEngine.ModSecurity} obsługuje serializowany JSON z
\code{modsec\_audit.log}. Parser:
\begin{enumerate}
    \item sprawdza poprawność JSON i strukturę transakcji,
    \item odczytuje \code{transaction\_id}, adres klienta i timestamp,
    \item parsuje request line oraz nagłówki,
    \item łączy fragmenty body zapisane w tablicy,
    \item odczytuje status odpowiedzi,
    \item zwraca \code{ModSecAuditEntry} albo szczegółowy błąd parsowania.
\end{enumerate}

Body z audit logu może zostać dołączone do obiektu \code{ApacheRequest} na
podstawie zgodnego \code{TransactionId}/\code{UniqueId}. Dzięki temu dane body
POST, których zwykle nie ma w \code{access.log}, mogą zostać przekazane do
detektora.

\section{Backend C\#}

\subsection{Worker odczytujący log}

\code{ApacheLogWorker} jest usługą działającą w tle. Ścieżka pliku jest pobierana
z konfiguracji \code{ApacheLog:Path}. Worker zapamiętuje pozycję w pliku, otwiera
plik z możliwością współdzielenia i odczytuje nowe linie. Błędy parsowania są
logowane, natomiast poprawne żądania trafiają do \code{LogEventProcessor}.

Drugi worker, \code{FlaggedEventRetentionWorker}, wykonuje okresowe czyszczenie
zdarzeń starszych niż 30 dni.

\subsection{LogEventProcessor}

Procesor wywołuje interfejs \code{IAnomalyDetector}. Jeżeli odpowiedź JSON zawiera
pole \code{is\_anomaly} ustawione na \code{true}, przekazuje żądanie, surową linię
logu i pełną odpowiedź detektora do repozytorium zdarzeń. Błędy komunikacji HTTP,
JSON i nieprawidłowego stanu są rejestrowane przez logger.

\subsection{SQLite i deduplikacja}

Repozytorium \code{SqliteFlaggedEventRepository} tworzy tabelę
\code{FlaggedEvents} przy uruchomieniu aplikacji. Przechowywane są między innymi:
\begin{itemize}
    \item czas zdarzenia i czas detekcji,
    \item surowa linia Apache,
    \item serializowane żądanie i pełny JSON detekcji,
    \item severity, metoda, URL i adres IP klienta.
\end{itemize}

Klucz zdarzenia jest obliczany jako SHA-256 surowej linii logu. Ogranicza to
wielokrotne zapisanie tego samego wpisu przy ponownym odczycie.

\section{Komunikacja z modułem detekcji Python}

\subsection{Kontrakt HTTP}

Backend korzysta z klienta \code{HttpClient}. Bazowy adres jest pobierany z
\code{PythonDetection:BaseUrl}, domyślnie jest to \code{http://127.0.0.1:8000/}.
Dla pojedynczego żądania adapter wysyła \code{POST /predict} z JSON-em:

\begin{lstlisting}
{
  "method": "POST",
  "url": "/login?next=/admin",
  "user_agent": "curl/8.0",
  "body": "username=admin&password=..."
}
\end{lstlisting}

Pole \code{body} ma wartość \code{request.Body.Content}, a przy braku body pusty
string. Dzięki temu body odczytane z ModSecurity jest analizowane przez Python.

\subsection{Odpowiedź detektora}

Endpoint \code{/predict} zwraca między innymi:
\begin{itemize}
    \item \code{anomaly\_score}, \code{threshold} i \code{score\_ratio},
    \item \code{is\_anomaly} oraz \code{severity},
    \item listę uruchomionych reguł i powodów,
    \item wyjaśnienia najważniejszych cech modelu.
\end{itemize}

Endpoint \code{/predict/batch} obsługuje listę do 1000 żądań i jest przeznaczony
do monitorowania w trybie wsadowym. Backend C\# używa obecnie endpointu
pojedynczego \code{/predict}.

\section{Poglądowe działanie detektora}

Moduł Python przyjmuje już znormalizowane dane, a nie surową linię
\code{access.log}. Dla każdego żądania wykonuje następujące kroki:

\begin{enumerate}
    \item normalizuje metodę, URL, User-Agent i body,
    \item rozdziela ścieżkę i parametry zapytania,
    \item wyznacza cechy długości, głębokości ścieżki, liczby parametrów,
    udziału znaków specjalnych i zakodowanych fragmentów,
    \item analizuje body, w tym liczbę parametrów i znaki specjalne,
    \item sprawdza sygnatury reguł bezpieczeństwa, takie jak SQL injection,
    XSS, traversal, ujawnienie plików lub command injection,
    \item przekazuje wektor cech do preprocessora i autoenkodera,
    \item łączy wynik modelu statystycznego z regułami deterministycznymi,
    \item buduje decyzję, poziom severity i objaśnienie XAI.
\end{enumerate}

Poglądowy schemat detekcji:

\begin{verbatim}
method, url, user_agent, body
              |
              v
      ekstrakcja cech i dekodowanie URL
              |
       +------+------+
       |             |
       v             v
   reguły       preprocessor
 bezpieczeństwa       |
       |              v
       |          autoencoder
       |              |
       +-------> decyzja i score
                       |
                       v
       is_anomaly, severity, reasons, top_features
\end{verbatim}

Schemat ma charakter poglądowy. Szczegółowe progi, sygnatury i cechy są
zdefiniowane w kodzie modułu \code{src/anomaly\_detection} oraz w plikach modeli.

\section{API i funkcjonalności aplikacji}

\subsection{Uwierzytelnianie}

Backend korzysta z uwierzytelniania cookie. Dostępne są endpointy:
\begin{itemize}
    \item \code{POST /api/auth/login} -- logowanie administratora,
    \item \code{POST /api/auth/logout} -- zakończenie sesji,
    \item \code{GET /api/auth/me} -- sprawdzenie bieżącej sesji.
\end{itemize}

Cookie sesji jest oznaczone jako HttpOnly i SameSite Strict. Endpointy zdarzeń
wymagają roli \code{Admin}.

\subsection{Zdarzenia anomalii}

\begin{itemize}
    \item \code{GET /api/anomalies} -- lista zdarzeń z paginacją i filtrem severity,
    \item \code{GET /api/anomalies/\{id\}} -- szczegóły zdarzenia,
    \item \code{GET /} -- prosty status backendu.
\end{itemize}

Szczegóły zawierają surowy wpis, sparsowane żądanie, wynik detektora i dane
wyjaśniające. Wyświetlane treści pochodzące z żądania są traktowane jako dane
niezaufane.

\subsection{Funkcjonalności frontendu}

Frontend React/TypeScript zapewnia:
\begin{enumerate}
    \item formularz logowania administratora,
    \item pobieranie listy oflagowanych żądań,
    \item prezentację adresu klienta, czasu, metody, URL i severity,
    \item otwieranie szczegółów zdarzenia,
    \item prezentację surowej linii Apache, żądania i odpowiedzi detektora,
    \item wylogowanie i zamknięcie widoku szczegółów.
\end{enumerate}

Aplikacja frontendowa korzysta z Vite i przekazuje żądania \code{/api} do backendu
podczas pracy deweloperskiej.

\section{Stos technologiczny}

\begin{longtable}{p{0.25\textwidth}p{0.62\textwidth}}
    \caption{Technologie użyte w projekcie}\label{tab:stack}\\
    \toprule
    \textbf{Obszar} & \textbf{Technologie} \\
    \midrule
    Backend & C\#, .NET 10, ASP.NET Core, kontrolery MVC, HttpClient \\
    Przechowywanie & SQLite, Microsoft.Data.Sqlite \\
    Parsery & C\#, Apache access.log, ModSecurity Serial JSON \\
    Detektor & Python, FastAPI, Uvicorn, PyTorch, pandas, NumPy,
    scikit-learn, joblib \\
    Frontend & React 19, TypeScript, Vite, CSS \\
    Testy & xUnit/.NET test runner, pytest, testy TypeScript przez kompilację
    projektu frontendowego \\
    Dokumentacja & Markdown i LaTeX \\
    \bottomrule
\end{longtable}

\section{Założenia techniczne i niefunkcjonalne}

\begin{itemize}
    \item System pracuje asynchronicznie: worker odczytuje log, a klient HTTP
    komunikuje się z detektorem z użyciem \code{CancellationToken}.
    \item Parser nie interpretuje body z samego access logu; body może zostać
    dołączone z ModSecurity.
    \item Surowe żądania i odpowiedzi detektora są przechowywane wyłącznie dla
    zdarzeń oznaczonych jako anomalne.
    \item Dane wejściowe żądań są potencjalnie kontrolowane przez atakującego i
    nie powinny być traktowane jako HTML ani wykonywany kod.
    \item Detektor może działać jako osobny proces FastAPI na tym samym hoście
    albo jako usługa dostępna pod adresem skonfigurowanym w backendzie.
    \item Model detekcji wymaga obecności plików modelu w katalogu
    \code{Deep Learning Detection Module/models}.
\end{itemize}

\section{Uruchomienie i wdrożenie}

Szczegółowe instrukcje uruchomienia znajdują się w pliku \code{UserManual.md}.
Główny wariant zakłada Linux i katalog główny repozytorium jako bieżący katalog.
Typowa kolejność jest następująca:

\begin{enumerate}
    \item utworzenie i aktywacja środowiska Python w katalogu detektora,
    \item uruchomienie FastAPI na porcie 8000,
    \item konfiguracja \code{ApacheLog:Path}, \code{Storage:DatabasePath}
    i \code{PythonDetection:BaseUrl},
    \item uruchomienie API ASP.NET Core,
    \item uruchomienie frontendu Vite.
\end{enumerate}

W środowisku produkcyjnym należy dodatkowo skonfigurować usługę systemową,
reverse proxy, TLS, kopie zapasowe SQLite i właściwe zarządzanie sekretami.

\section{Testowanie i weryfikacja}

Projekt zawiera testy parsera Apache, parsera ModSecurity, serializacji modeli,
repozytorium SQLite i klienta detektora. Weryfikacja wykonana dla aktualnego
stanu projektu wykazała:
\begin{itemize}
    \item 195 zaliczonych testów .NET,
    \item 36 zaliczonych testów Python,
    \item poprawny build frontendu przez \code{npm run build},
    \item poprawne uruchomienie endpointów \code{/health} i \code{/predict}
    z załadowanym modelem.
\end{itemize}

\todo{Uzupełnić diagramy architektury i sekwencji w formie grafik, jeśli dokument
będzie składany jako pełna praca projektowa.}

\section{Organizacja pracy}

\todo{Uzupełnić sposób planowania zadań, narzędzia komunikacji, zasady przeglądu
kodu, strategię branchowania, harmonogram i podział pracy w kolejnych tygodniach.}

\section{Uwarunkowania prawne i społeczne}

System przetwarza adresy IP, identyfikatory transakcji, nagłówki HTTP i fragmenty
żądań. W zależności od jurysdykcji dane te mogą stanowić dane osobowe lub dane
umożliwiające identyfikację użytkownika.

\todo{Uzupełnić analizę podstawy prawnej przetwarzania logów, okresu retencji,
zasad anonimizacji, kontroli dostępu, obowiązków administratora danych oraz
informowania użytkowników.}

\section{Analiza rynku i konkurencji}

System należy do kategorii narzędzi monitoringu bezpieczeństwa aplikacji,
analizy logów, SIEM oraz detekcji anomalii. Rozwiązanie łączy parser Apache,
reguły bezpieczeństwa i model uczenia maszynowego w lekkiej architekturze
lokalnej.

\todo{Uzupełnić analizę rynku, grupę docelową, porównanie z narzędziami SIEM,
WAF i systemami IDS/IPS, przewagi rozwiązania, koszty wdrożenia oraz model
utrzymania.}

\section{Ograniczenia i dalszy rozwój}

Aktualne ograniczenia obejmują zależność parsera od określonego rozszerzonego
formatu logu, konieczność dostępności modelu Python oraz brak rozbudowanej
orkiestracji wdrożenia produkcyjnego.

Potencjalne kierunki rozwoju:
\begin{itemize}
    \item bezpieczne buforowanie niepełnych linii podczas rotacji i zapisu logu,
    \item obsługa większej liczby formatów Apache i innych źródeł logów,
    \item wsadowe wysyłanie żądań do endpointu \code{/predict/batch},
    \item rozbudowane filtrowanie, sortowanie i paginacja w interfejsie,
    \item monitoring dostępności i czasu odpowiedzi detektora,
    \item wdrożenie z użyciem kontenerów i usług systemowych.
\end{itemize}

\section{Podsumowanie}

\project{} (System Ciągłej Hybrydowej Identyfikacji Zagrożeń i Anomalii) jest
systemem łączącym cykliczne monitorowanie logów Apache, deterministyczne reguły
bezpieczeństwa i analizę anomalii wykonywaną przez model uczenia maszynowego
w module Python. Nazwa odzwierciedla zastosowany sposób działania: odczyt
kolejnych wpisów logu, analizę żądań HTTP z wykorzystaniem podejścia
hybrydowego oraz identyfikowanie zdarzeń wskazujących na potencjalne zagrożenia
lub nietypowe zachowania.

Rozdzielenie backendu, parserów, detektora i frontendu ułatwia niezależne
testowanie oraz wymianę poszczególnych komponentów. Najważniejszym kontraktem
integracyjnym jest znormalizowane żądanie zawierające metodę, URL, User-Agent
i body, które backend wysyła do endpointu detektora.

\end{document}